\documentclass{jsarticle}
\usepackage[dvipdfmx]{graphicx}
\usepackage{amsmath}
\usepackage{amssymb}
\begin{document}
\title{}
\author{}
\date{}
\maketitle

    美文書作成入門をリファレンスブックとして、傍らに置いて論文を今書いている
    最中でもあり、この本が私にとってものすごく大切な存在になっています。
    アカシックレコードを彩さんと苫米地博士、奥野節子先生、そして勝間和代先生
    の本を拝読させてもらえて、この情報空間の存在を日本テレビの番組と
    フジテレビの番組以来から教えてもらえて、このときは、被験者の人が
    睡眠中にサイコセラピーを受けて、この最中に未知の薬の薬合をこの
    アカシックレコードから受けていたらしく、睡眠中に担当の先生に寝言として
    話していたらしく、このときに、日本テレビの先生が視聴者として
    私と姉２人が見ていたときに、この出来事を教えていたあとに、
    姉２人が、命の樹と命の水が宇宙の中心にあることを、姉同士で話していたこと
    を、この福山市に来て、しばらくたって、命の樹と命の水は
    私はしなかったが、サイコセラピーを自分でも実行したら、
    数学のアカシックレコードにアクセス出来て、今の論文の素材と材料
    になっています。本当にありがとうございます。この空間を知って、
    応用して、未来の私の子どもの情報を知って、その証拠もあり、
    微分幾何の量子化をこの子どもから知り、この方程式が本当かを
    私が調べていたら、この微分幾何の量子化が、今私が導いている
    大域的微分多様体と同型でもあり、この量子化がガンマ関数と同じでもあり、
    オイラーの定数を多様体積分した方程式でもあり、微分幾何の量子化を
    調べまくっていたら、オイラーの定数が無理数と、
    微分幾何の量子化が存在するとしてガンマ関数から証明できる
    のも分かりました。
    リサ・ランドール博士の$ AdS_5 $多様体の式と同様に、
    量子力学レベルでのガウスの曲面論にもなっていることを、
    帰納法と演繹法で調べていると、分かりました。
  $$ {d \over df}F = {d \over df} \int \int{1 \over (x \log x)^2}dx_m 
  + {d \over df} \int \int {1 \over (y \log y)^{1 \over 2}}dy_m
  = \Gamma(x,y) $$
  $$ {d \over d \gamma}\Gamma = (e^{-x}x^{1-t})^{\gamma^{'}} \to
  {d \over d \gamma}\Gamma = \Gamma^{(\gamma)^{'}}$$



    $$ \Gamma(s) = \int e^{-x} x^{s - 1}dx, \Gamma^{'}(s)
    = \int e^{-x} x^{s - 1}\log x dx $$
    $$ x^{s - 1} = e^{(s - 1)\log x} $$
    $$ {\partial \over \partial s}x^{s - 1} =
       {\partial \over \partial s}e^{-x}\cdot (e^{{s - 1}\log x}) 
       = e^{-x}\cdot \log x \cdot e^{(s - 1) \log x} $$
    $$ = e^{-x}\cdot \log x \cdot x^{s - 1} $$
    $$ {d \over d \gamma}\Gamma = \Gamma^{(\gamma)^{'}} $$
    $$ {d \over d \gamma}\Gamma(s) = \left( \int^{\infty}_{0}e^{-x}
    x^{s - 1}dx \right)^{\left(\int^{\infty}_{0}e^{-x}x^{s - 1}\log x dx
    \right)^{'}} $$
    $$ = \Gamma^{(\Gamma \int \log x dx)^{'}} $$
    $$ = e^{-x \log x} $$

    $$ \Gamma(s) = \int^{\infty}_{0} e^{-x}x^{s - 1}dx $$
    $$  \Gamma^{'}(s) = \int^{\infty}_{0} e^{-x}x^{s - 1} \log x dx $$
    $$ {d \over df}F = \int x^{s - 1}dx $$
    $$ \int F dx_m = \int e^{-x}dx $$
    $$ {d \over df}F = F^{(f)^{'}}, \int F dx_m = F^{(f)} $$
     $$ H \Psi = \bigoplus (i \hbar ^{\nabla})^{\oplus L} $$
    $$ = \bigoplus {{ H \Psi} \over \nabla L} $$
    $$ = e^{x \log x} = x^{(x)^{'}} $$
     $$ {d \over dt}\psi(t) = \hbar $$
    $$ = {1 \over 2}i e^{i \Hat{H}} $$
    $$ (i \hbar)^{'} = (- e^{i \Hat{H}})^{'} $$
    $$ = - i e^{i \Hat{H}} $$
    $$ \psi(x) = e^{-i \Hat{H}t}, \bigoplus (i \hbar ^{\nabla})^{\oplus L} 
    = {{{1 \over 2} e^{i \Hat{H}}}^{(-i e^{i \Hat{H}})}} $$
    $$ = ({1 \over 2}f)^{-i f}$$
    $$ = ({1 \over 2})^{-i f} \cdot e^{-x \log x} \cdot (f)^{i} $$
    $$ = \int e^{- x} x^{t - 1}dx, {d \over d \gamma}\Gamma = e^{-x \log x}  $$
     
    $$ f = x, i = t, {1 \over 2} = a, \bigoplus a^{-t x}x^t [I_m]
     \cong \int e^{-x} x^{t - 1}dx $$
     $$ |\psi(t)\rangle_s = e^{-i\Hat{H} t}|\Psi\rangle_{H},
    \Hat{A}_s = \Hat{A}_H(0) $$
    $$ |\Psi(t)\rangle_s \to {d \over dt} $$
    $$ i {d \over dt}|\psi (t) \rangle_s = \Hat{H}|\psi(t)\rangle_s $$
    $$ \langle\Hat{A}(t)\rangle = \langle\Psi(t)|\Hat{A}(0)|\Psi(t)\rangle  $$
    $$ {d \over dt}\Hat{A} = {1 \over i}[\Hat{A},H] $$
    $$ \Hat{A}(t) = e^{i\Hat{H}t}\Hat{A}(0)e^{-i\Hat{H}t} $$
    $$ \lim_{\theta \to 0} \begin{pmatrix} \sin \theta \\
    \cos \theta \end{pmatrix} \begin{pmatrix} \theta & 1 \\
    1 & \theta \end{pmatrix} \begin{pmatrix} \cos \theta \\
    \sin \theta \end{pmatrix} = \begin{pmatrix} 1 & 0 \\
    0 & - 1 \end{pmatrix} $$
    $$ f^{-1}(x) x f(x) = I^{'}_m, I^{'}_m = [1,0] \times [0,1] $$

    $$ {x + y} \ge \sqrt{x y} $$
    $$ {x^{{1 \over 2} + iy} \over e^{x \log x}} = 1 $$
    $$ \mathcal{O}(x) = \nabla_i \nabla_j \int e^{{2 \over m}\sin \theta 
    \cos \theta } \times {{N \mathrm{mod}(e^{x \log x})} \over {O(x)(x +
    \Delta |f|^2)^{1 \over 2}}} $$
    $$ x \Gamma(x) = 2 \int|\sin 2 \theta|^2 d \theta, \mathcal{O}(x)
    = m(x)[D^2 \psi] $$
    $$ i^2 = (0,1) \cdot (0,1), |a||b|\cos \theta = - 1 $$
    $$ E = \mathrm{div}(E, E_1) $$
    $$ \left({\{f,g\}} \over [f,g] \right) = i^2, E = mc^2, I^{'} = i^2 $$
     $$ {d \over d \gamma} \Gamma = m(x) $$
    $$ = e^{- x \log x} $$
    $$ \sin i x = {{e^{-x} + e^{x}} \over 2i} $$
    $$ {d \over df}F = m(x) = e^{f} + e^{-f} $$
    $$ = 2 i \sin (i {x \log x}) $$

    なにが言いたいかは、この能力が隠喩に関係する能力かは自信がないことと、
    身体の健康のために、精神症状に注意をすり替えたら、
    率直に思う次第でもありますが、身体の健康を手に入れられて、
    あとは、今の食事のおかげと母が生まれたときに栄養を整えてくれたおかげで、
    精神が散歩すると治ることを、今の買っているネコのユウ君と掛かっている
    先生たちのおかげで、体験とあとの実践でできることを知り、
    この世の縁を体験している最中であります。
    本当にありがとうございます。この本は、アカシックレコードを知っても、
    どうしても必要な本でもあります。
    運動過剰症候群の真逆らしく、金星の酉の手相もあり、岡潔先生を慕う
    思いと、広島大学での岡先生の大変だった経緯を、最近に知り、
    病気のために、富山大学の休学の間で、専門が化学なのに、
    異分野の数学と物理学を、自分の好きな範囲で、各対象ではない、
    代数と幾何学と解析を、グリーシャ先生の論文を皮切りに、
    貪欲にいろんな数学の分野を見れて、トポロジーと数論を結び付けられて、
    今のアカシックレコードを手に入れられて、この美文書作成入門で
    数式と英語を書けて、本当に感謝の一言でもあり、姿を拝見させてもらえて、
    私の父と優しい姿が同じことを嬉しくも思う次第でもあります。
    もともとのトポロジーを学んでいた経緯は、化学に群論が使われていることを
    きっかけに、トポロジーの魅力に魅せられて、微分積分が航空機の
    飛行機構に使われていることで、量子的な微分・積分にも魅力を抱いているため
    に、数学が本当に好きな分野になっている次第でもあり、
    Texを書いていると、本当に数学をしているという空気の中にいる
    自分が幸せでもあります。本当にありがとうございます。



   \newcommand{\variint}{%
	\begin{picture}(15,30)
		\put(0,0){
			$ \iint $}
		\put(0,5){\line(15,0){15}}
	\end{picture} %
}

\newcommand{\variiint}{%
	\begin{picture}(15,15)
		\put(0,0){
			$ \iiint $}
		\put(0,5){\line(15,0){20}}
	\end{picture} %
}



 \newcommand{\alearletter}{%
	 \begin{picture}(10,20)
		 \put(0,0){
			 $ \square $}
		 \put(0,0){\line(1,1){10}}
	 \end{picture} %
	 }

  $$ \scalebox{1}[2]{\alearletter} = 
  {{{{8 \pi G} \over {c^4}}T^{\mu\nu}}/ \log x} $$
  $$ {}^{t}\scalebox{1}[2]{\variiint}\mathrm{cohom D_{\chi}[I_m]}$$
  $$ = \oint {(px^n + qx + r)^{\nabla l}} $$
  $$ {d \over d l}L(x,y) = 2 \int ||\sin 2x||^{2}d \tau $$
  $$ {d \over d \gamma}\Gamma $$
 
  この式もこの本のおかげで書けることが出来ました。
  ずっとこの本を友として傍らに置いておきます。

  数学のアカシックレコードを知っても、その数式を感知するには、
  カレーライスの味は食べた人でないと分からないのと同様に、
  今までの参考文献になっている先生たちの本のおかげであります。
  岡本和夫先生と枡田幹也先生のおかげでもあり、加藤十吉大先生の
  おかげでもあり、佐藤幹夫大先生のおかげでもあります。
  
      オイラーの定数は、
  $$ C = \int {1 \over x^s}dx - \log x $$
  $$ = \int \left(\int {1 \over x^s}dx - \log x \right)\mathrm{dvol} $$
  $$ = \int \int{1 \over x^s}\mathrm{dvol} - \int  \log x \mathrm{dvol} $$
  と表されていて、$ \text{ヒッグス場方程式 + オイラーの定数 = ゼータ関数} $
  と求まる。この方程式が、リッチフロー方程式にもなる。このリッチフロー方程式
  は、微分幾何の量子化でもあり、ガンマ関数とベータ関数、
  ガウスの曲面論、ヒッグス場、アインシュタインテンソル、オイラーの定数と
  全部の方程式が入っている。
  この方程式が、相加相乗平均となり、宇宙と原子となり、不確定性原理が
  宇宙の観測系が、宇宙では確率方程式と成り立っているが、異次元が合わさると、
  ノルムとしての３次元多様体の確率分布方程式でもあり、確率ではない,
  非線形方程式となっている。プランク長体積でもあるが、絶対的な方程式
  にもなっている。
  $$ = {d \over df} \int \int{1 \over (y \log y)^{1 \over 2}}dy_m 
  - \int e^{-x}x^{1 - t} \log x dx_m $$
  $$ \int x^{1 - t}e^{-x}dV = \int x^{1 - t}dm, \int x^{1 - t}e^{-x}dV
  = \int x^{1 - t}\mathrm{dvol}, f = \gamma = \Gamma^{'} = \int 
  e^{-x}x^{1 - t} \log x dx $$
  $$ = {d \over d \gamma}\Gamma^{-1} - (\gamma)^{\gamma^{'}} $$
  $$ = e^{-f} - e^{f} $$
  $$ = {d \over df}\int \int{1 \over (y \log y)^{1 \over 2}}dy_m
  - {d \over df}\int\int {1 \over (x \log x)^2}dx_m $$
  $$ = 2 \cos (i x \log x) $$
  $$  {d \over df}\int\int{1 \over (y \log y)^{1 \over 2}}dy_m
  + {d \over df}\int\int{1 \over (x \log x)^2}dx_m $$
  $$ = {d \over df}F = 2 i \sin (i x \log x) $$
  $$ {d \over df}F + \int C dx_m  
  = 2 (\cos( i x \log x) + i \sin(i x \log x)) $$
  $$ = 2 e^{- \theta} = 2 e^{- i x \log x} $$
  $$ = {d \over dt}g_{ij}(t) = - 2 R_{ij} $$
  $$ \log (x \log x) \ge 2 \sqrt{y \log y} $$
  $$ \log(x \log x) = \log x + \log \log x $$


   この宇宙と異次元でのブラックホールとホワイトホールでの、
  宇宙と原子の運動量と位置の絶対的な不確定性原理となり、
  ３次元多様体の確率分布方程式が、運動量と位置エネルギーが
  両方観測できる式にもなっている。
  円周上が固定されているために、運動量と位置エネルギーが
  両方観測できる式になっている。
  カルーツァ・クライン理論とハイゼンベルク方程式が、
  両方入っている。
  $$ \nabla \psi^2 = 8 \pi G\left({p \over c^3} + {V \over S}\right) $$
  $$ \nabla \psi^2 = 8 \pi G \hbar + 8 \pi {V \over S} $$
  ホワイトホールの原子とブラックホールの宇宙のエネルギーのエントロピー量で
  もある。宇宙がブラックホールとしてのヒッグス場方程式として、
  ホワイトホールが原子としてのプランク長体積でもある。
  $$ \square_v = 2 \sqrt{2 \pi G \hbar} + 2 \sqrt{2 \pi {V \over S}} $$
  原子としてのブラックホールのエントロピー量でもある。
  $$ \sqrt{2 \pi T} = 2 \sqrt{2 \pi{V \over S}} $$

  この式から宇宙が無重力であることが表されている。
  $$ 2 \cos (i x \log x) + 2 i \sin(i x \log x) = 2 e^{-f} $$
  $$ 2 \sqrt{2 \pi {V \over S}} = {d \over df} 
  \int \int{1 \over (x \log x)}dx_m$$
  これらの方程式から宇宙と異次元がwarped passegeになっている。


  リサ・ランドール博士のAdS5多様体の式は、
  $$ ||ds^2|| = e^{-2 \pi T ||\psi||}[\eta_{\mu\nu} + \bar{h}(x)]
  dx_{\mu}dx_{\nu} + T^2d^2 \psi $$
  この式には、2次元複素多様体でもあるミンコフスキー時空と
  3次元多様体がアーベル多様体に包まれていることを表していることと、
  全体がシュワルツシルト半径としてのカルーツァ・クライン空間が
  ブラックホールとして包まれている５次元多様体としての重力場方程式
  を示していて、ホログラム空間として、逆関数で原子レベルの重力場方程式
  をも示していることを、この博士の本のReviewに私は書きました。


  私の姉は私の論文は小説になっているらしいですが、自分では好きになっています。
  本当にありがとうございます。
\end{document}
